\documentclass[11pt]{article}
\usepackage{amsmath, amssymb, geometry, xcolor}
\geometry{letterpaper, margin=1in}

\title{\textbf{Matrix Operations: Comprehensive Problem Set}}
\author{Curated Practice based on Homework 2.1}
\date{}

\begin{document}

\maketitle
\tableofcontents
\newpage

% ==========================================
% SECTION 1: CONCEPT CHECKLIST & MAPPING
% ==========================================
\section{Concept Checklist \& Question Mapping}
This checklist outlines all the core concepts, problem types, and mathematical tricks extracted from the provided text. The numbers denote which problems in the set evaluate that specific concept.

\begin{itemize}
    \item \textbf{C1: Mechanical Arithmetic \& Dimensional Analysis}
    \begin{itemize}
        \item Matrix Addition/Subtraction and Scalar Multiplication \textit{(Questions: 1, 2)}
        \item Verifying Legal Operations via Matrix Dimensions \textit{(Questions: 3, 4, 5)}
    \end{itemize}
    
    \item \textbf{C2: Matrix Multiplication}
    \begin{itemize}
        \item The Row-Column Dot Product Rule \textit{(Questions: 6, 7)}
        \item Distributive Law Expansion \textit{(Questions: 14, 15)}
    \end{itemize}
    
    \item \textbf{C3: Parameterized Commutativity (Harder)}
    \begin{itemize}
        \item Solving for variables forcing $AB = BA$ \textit{(Questions: 8, 9)}
        \item Proving non-commutativity exists (No Solution systems) \textit{(Question: 10)}
        \item Structural constraints for commutative matrices \textit{(Question: 11)}
    \end{itemize}

    \item \textbf{C4: Null Space \& Zero Divisor Construction (Harder)}
    \begin{itemize}
        \item Creating non-zero matrices $B$ such that $AB = 0$ \textit{(Questions: 12, 13)}
        \item Nilpotent Matrices ($A^k = 0$ despite $A \neq 0$) \textit{(Questions: 29, 30)}
    \end{itemize}

    \item \textbf{C5: Inner and Outer Products}
    \begin{itemize}
        \item Scalar output of Inner Products ($u^T v$) \textit{(Question: 16)}
        \item Matrix output of Outer Products ($u v^T$) \textit{(Question: 17)}
        \item Quadratic Forms ($x^T A x$) \textit{(Question: 18)}
    \end{itemize}

    \item \textbf{C6: The Trace Operator}
    \begin{itemize}
        \item Sum of main diagonals \textit{(Question: 19)}
        \item Trace Commutativity Proofs ($\text{tr}(AB) = \text{tr}(BA)$) \textit{(Questions: 20, 23)}
    \end{itemize}

    \item \textbf{C7: Transpose Properties}
    \begin{itemize}
        \item Transpose of a Product ($(AB)^T = B^T A^T$) \textit{(Question: 21)}
        \item Proving symmetry in matrix combinations \textit{(Question: 22)}
    \end{itemize}

    \item \textbf{C8: Diagnostic "Find the Flaw" Testing (Harder)}
    \begin{itemize}
        \item Matrix Expansion Pitfalls (Assuming Commutativity) \textit{(Question: 24)}
        \item The Zero-Product Fallacy \textit{(Question: 25)}
        \item Outer Product Invertibility Fallacy \textit{(Question: 26)}
        \item Transpose Dimension Trap \textit{(Question: 27)}
    \end{itemize}

    \item \textbf{C9: Real-World Applications \& Modeling}
    \begin{itemize}
        \item Financial Portfolio Variance ($w^T \Sigma w$) \textit{(Question: 28)}
        \item Markov Chains and Stochastic Transitions \textit{(Question: 29)}
    \end{itemize}
\end{itemize}

\newpage
% ==========================================
% SECTION 2: PROBLEM SETS
% ==========================================
\section{Problem Sets}

\textbf{Part 1: Mechanical Arithmetic \& Dimensional Analysis}
\vspace{0.3cm}

\noindent \textbf{Problem 1.} Let $A = \begin{bmatrix} 1 & -2 \\ 3 & 0 \end{bmatrix}$ and $B = \begin{bmatrix} -1 & 4 \\ 2 & -1 \end{bmatrix}$. Compute $2A - 3B$.

\noindent \textbf{Problem 2.} Let $X = \begin{bmatrix} 4 & 5 \\ -1 & 2 \end{bmatrix}$ and $Y = \begin{bmatrix} 0 & 3 \\ 1 & -1 \end{bmatrix}$. Compute $X + 2Y^T$.

\noindent \textbf{Problem 3.} Matrix $A$ is $4 \times 3$, Matrix $B$ is $3 \times 5$, and Matrix $C$ is $5 \times 2$. Is the product $ABC$ defined? If so, what are its dimensions?

\noindent \textbf{Problem 4.} Matrix $U$ is $3 \times 2$, Matrix $V$ is $2 \times 4$, and Matrix $W$ is $3 \times 4$. Determine if the operation $U V + W$ is legal. If so, what are the final dimensions?

\noindent \textbf{Problem 5.} Let $C = \begin{bmatrix} 1 & 2 \\ 3 & 4 \\ 5 & 6 \end{bmatrix}$ and $D = \begin{bmatrix} 1 & 0 & -1 \\ -2 & 1 & 0 \end{bmatrix}$. Compute $CD$. Is $DC$ defined? If so, what are its dimensions?

\vspace{0.5cm}
\textbf{Part 2: Matrix Multiplication}
\vspace{0.3cm}

\noindent \textbf{Problem 6.} Let $A = \begin{bmatrix} 1 & -1 \\ 2 & 0 \\ 3 & 2 \end{bmatrix}$ and $x = \begin{bmatrix} 4 \\ -2 \end{bmatrix}$. Compute $Ax$.

\noindent \textbf{Problem 7.} Let $A = \begin{bmatrix} 2 & 1 \\ -1 & 3 \end{bmatrix}$ and $B = \begin{bmatrix} 0 & -2 \\ 4 & 1 \end{bmatrix}$. Compute $AB$ and $BA$.

\vspace{0.5cm}
\textbf{Part 3: Parameterized Commutativity}
\vspace{0.3cm}

\noindent \textbf{Problem 8.} Let $A = \begin{bmatrix} 2 & -1 \\ 3 & 1 \end{bmatrix}$ and $B = \begin{bmatrix} 1 & x \\ y & 2 \end{bmatrix}$. Find the specific values of $x$ and $y$ such that $AB = BA$.

\noindent \textbf{Problem 9.} Let $A = \begin{bmatrix} 1 & 1 \\ 0 & 1 \end{bmatrix}$ and $B = \begin{bmatrix} a & b \\ c & d \end{bmatrix}$. What conditions must be placed on the variables $a, b, c,$ and $d$ to guarantee $AB = BA$?

\noindent \textbf{Problem 10 (Trick Question).} Let $A = \begin{bmatrix} 4 & -2 \\ -2 & 1 \end{bmatrix}$ and $B = \begin{bmatrix} 2 & k \\ -1 & 3 \end{bmatrix}$. Find the value of $k$ such that $AB = BA$. If none exists, explain why.

\noindent \textbf{Problem 11.} Let $D = \begin{bmatrix} 2 & 0 \\ 0 & 3 \end{bmatrix}$. Find all $2 \times 2$ matrices $M = \begin{bmatrix} a & b \\ c & d \end{bmatrix}$ such that $DM = MD$. Describe the structure of $M$.

\vspace{0.5cm}
\textbf{Part 4: Null Space \& Zero Divisors}
\vspace{0.3cm}

\noindent \textbf{Problem 12.} Let $A = \begin{bmatrix} 3 & 6 \\ 1 & 2 \end{bmatrix}$. Construct a $2 \times 2$ matrix $B$, where $B \neq 0$, such that $AB = \begin{bmatrix} 0 & 0 \\ 0 & 0 \end{bmatrix}$.

\noindent \textbf{Problem 13.} Let $C = \begin{bmatrix} 1 & -3 \\ -2 & 6 \end{bmatrix}$. Construct a non-zero $2 \times 2$ matrix $D$ such that $DC = 0$. (Notice the left-multiplication).

\vspace{0.5cm}
\textbf{Part 5: Matrix Algebra \& Expansions}
\vspace{0.3cm}

\noindent \textbf{Problem 14.} Expand and simplify the expression $(A - 2B)(A + B)$. Assume $A$ and $B$ are $n \times n$ matrices that do \textbf{not} commute.

\noindent \textbf{Problem 15.} Simplify the expression $(A+I)^3$, where $I$ is the Identity matrix. Justify why the standard binomial expansion holds in this specific case.

\vspace{0.5cm}
\textbf{Part 6: Inner Products, Outer Products, and Trace}
\vspace{0.3cm}

\noindent \textbf{Problem 16.} Let $u = \begin{bmatrix} 1 \\ -2 \\ 3 \end{bmatrix}$ and $v = \begin{bmatrix} 4 \\ 0 \\ -1 \end{bmatrix}$. Compute the inner product $u^T v$.

\noindent \textbf{Problem 17.} Using the same vectors $u$ and $v$ from Problem 16, compute the outer product $uv^T$.

\noindent \textbf{Problem 18.} Let $x = \begin{bmatrix} 1 \\ 2 \end{bmatrix}$ and $A = \begin{bmatrix} 3 & -1 \\ 0 & 4 \end{bmatrix}$. Compute the quadratic form $x^T A x$. What are the dimensions of the final answer?

\noindent \textbf{Problem 19.} Let $M = \begin{bmatrix} 5 & 2 & -1 \\ 4 & -3 & 6 \\ 7 & 0 & 9 \end{bmatrix}$. Compute $\text{tr}(M)$.

\noindent \textbf{Problem 20.} Let $A = \begin{bmatrix} 2 & -1 \\ 0 & 3 \end{bmatrix}$ and $B = \begin{bmatrix} 1 & 4 \\ 5 & -2 \end{bmatrix}$. Compute $AB$ and $BA$. Verify that $\text{tr}(AB) = \text{tr}(BA)$.

\vspace{0.5cm}
\textbf{Part 7: Transpose Properties \& Proofs}
\vspace{0.3cm}

\noindent \textbf{Problem 21.} Let $A = \begin{bmatrix} 1 & 0 \\ -1 & 2 \end{bmatrix}$ and $B = \begin{bmatrix} 3 & 1 \\ 4 & -1 \end{bmatrix}$. Compute $(AB)^T$ and verify it is exactly equal to $B^T A^T$.

\noindent \textbf{Problem 22.} Prove algebraically that for any arbitrary matrix $A$ of dimension $m \times n$, the resulting matrix of $A^T A$ is always square and symmetric. 

\noindent \textbf{Problem 23.} Let $u = \begin{bmatrix} u_1 \\ u_2 \end{bmatrix}$ and $v = \begin{bmatrix} v_1 \\ v_2 \end{bmatrix}$. Show algebraically that the trace of their outer product equals their inner product. That is, prove $\text{tr}(uv^T) = u^T v$.

\vspace{0.5cm}
\textbf{Part 8: Find the Flaw (Diagnostic)}
\vspace{0.3cm}

\noindent \textbf{Problem 24.} \textit{Prompt:} Expand $(A + B)^2$.\\
\textit{Flawed Solution:} $(A+B)^2 = (A+B)(A+B) = A^2 + 2AB + B^2$.\\
\textit{Task:} Identify the flaw in reasoning and provide the correct expansion.

\noindent \textbf{Problem 25.} \textit{Prompt:} Solve for $X$ in the equation $AX = 0$, where $A = \begin{bmatrix} 2 & 4 \\ 1 & 2 \end{bmatrix}$ and $X = \begin{bmatrix} x_1 \\ x_2 \end{bmatrix}$.\\
\textit{Flawed Solution:} By the zero-product property, if $AX = 0$ and $A \neq 0$, then $X$ must be the zero vector.\\
\textit{Task:} Identify the flaw and provide the correct solution set for $X$.

\noindent \textbf{Problem 26.} \textit{Prompt:} Let $u$ and $v$ be non-zero $3 \times 1$ vectors. Compute $(uv^T)^{-1}$.\\
\textit{Flawed Solution:} Using the inverse properties, $(uv^T)^{-1} = (v^T)^{-1} u^{-1}$.\\
\textit{Task:} Identify the two severe mathematical flaws in this solution regarding vectors and outer products.

\noindent \textbf{Problem 27.} \textit{Prompt:} Let $A$ be $3 \times 4$ and $B$ be $4 \times 5$. Determine the dimensions of $(AB)^T$.\\
\textit{Flawed Solution:} By properties of transposes, $(AB)^T = A^T B^T$. Thus, multiplying a $4 \times 3$ by a $5 \times 4$ matrix yields a $5 \times 3$ matrix.\\
\textit{Task:} Identify the flaw and state the correct property and final dimensions.

\vspace{0.5cm}
\textbf{Part 9: Application and Abstract Models}
\vspace{0.3cm}

\noindent \textbf{Problem 28 (Financial Modeling).} A portfolio consists of two assets. The weight vector is $w = \begin{bmatrix} 0.5 \\ 0.5 \end{bmatrix}$ and the covariance matrix is $\Sigma = \begin{bmatrix} 0.16 & 0.02 \\ 0.02 & 0.25 \end{bmatrix}$. Compute the portfolio variance $\sigma_p^2 = w^T \Sigma w$.

\noindent \textbf{Problem 29 (Markov Chains).} Let $P = \begin{bmatrix} 0.9 & 0.2 \\ 0.1 & 0.8 \end{bmatrix}$ be a column-stochastic transition matrix for customer retention. Vector $v_0 = \begin{bmatrix} 1 \\ 0 \end{bmatrix}$ means $100\%$ of users start on Product A. Compute $v_2 = P^2 v_0$, representing the distribution after 2 cycles.

\noindent \textbf{Problem 30 (Nilpotent Structure).} Let $A = \begin{bmatrix} 0 & 1 & 0 \\ 0 & 0 & 1 \\ 0 & 0 & 0 \end{bmatrix}$. Compute $A^2$ and $A^3$. What concept from the text does this heavily reinforce regarding zero-products in matrices?

\newpage
% ==========================================
% SECTION 3: DETAILED SOLUTIONS
% ==========================================
\section{Detailed Solutions}

\textbf{Solution 1.} 
$2A = \begin{bmatrix} 2 & -4 \\ 6 & 0 \end{bmatrix}$, $3B = \begin{bmatrix} -3 & 12 \\ 6 & -3 \end{bmatrix}$.\\
$2A - 3B = \begin{bmatrix} 2 - (-3) & -4 - 12 \\ 6 - 6 & 0 - (-3) \end{bmatrix} = \begin{bmatrix} 5 & -16 \\ 0 & 3 \end{bmatrix}$.

\vspace{0.2cm}\noindent\textbf{Solution 2.}
$Y^T = \begin{bmatrix} 0 & 1 \\ 3 & -1 \end{bmatrix}$. $2Y^T = \begin{bmatrix} 0 & 2 \\ 6 & -2 \end{bmatrix}$.\\
$X + 2Y^T = \begin{bmatrix} 4 & 5 \\ -1 & 2 \end{bmatrix} + \begin{bmatrix} 0 & 2 \\ 6 & -2 \end{bmatrix} = \begin{bmatrix} 4 & 7 \\ 5 & 0 \end{bmatrix}$.

\vspace{0.2cm}\noindent\textbf{Solution 3.}
Using the size check cheat: $(4 \times 3) \times (3 \times 5) \rightarrow (4 \times 5)$. The inner 3s cancel. 
Next, $(4 \times 5) \times (5 \times 2) \rightarrow (4 \times 2)$. 
Yes, it is defined. Dimensions are \textbf{$4 \times 2$}.

\vspace{0.2cm}\noindent\textbf{Solution 4.}
$U$ is $3 \times 2$ and $V$ is $2 \times 4$. $UV$ yields a $3 \times 4$ matrix. 
$W$ is $3 \times 4$. Since $UV$ and $W$ have identical dimensions, they can be legally added. 
Final dimensions: \textbf{$3 \times 4$}.

\vspace{0.2cm}\noindent\textbf{Solution 5.}
$CD = \begin{bmatrix} 1 & 2 \\ 3 & 4 \\ 5 & 6 \end{bmatrix} \begin{bmatrix} 1 & 0 & -1 \\ -2 & 1 & 0 \end{bmatrix} = \begin{bmatrix} (1-4) & (0+2) & (-1+0) \\ (3-8) & (0+4) & (-3+0) \\ (5-12) & (0+6) & (-5+0) \end{bmatrix} = \begin{bmatrix} -3 & 2 & -1 \\ -5 & 4 & -3 \\ -7 & 6 & -5 \end{bmatrix}$.\\
$DC$ is $(2 \times 3) \times (3 \times 2)$, which is \textbf{defined} and yields a \textbf{$2 \times 2$} matrix.

\vspace{0.2cm}\noindent\textbf{Solution 6.}
$Ax = \begin{bmatrix} 1 & -1 \\ 2 & 0 \\ 3 & 2 \end{bmatrix} \begin{bmatrix} 4 \\ -2 \end{bmatrix} = \begin{bmatrix} 1(4) + (-1)(-2) \\ 2(4) + 0 \\ 3(4) + 2(-2) \end{bmatrix} = \begin{bmatrix} 6 \\ 8 \\ 8 \end{bmatrix}$.

\vspace{0.2cm}\noindent\textbf{Solution 7.}
$AB = \begin{bmatrix} 2 & 1 \\ -1 & 3 \end{bmatrix} \begin{bmatrix} 0 & -2 \\ 4 & 1 \end{bmatrix} = \begin{bmatrix} 4 & -3 \\ 12 & 5 \end{bmatrix}$.\\
$BA = \begin{bmatrix} 0 & -2 \\ 4 & 1 \end{bmatrix} \begin{bmatrix} 2 & 1 \\ -1 & 3 \end{bmatrix} = \begin{bmatrix} 2 & -6 \\ 7 & 7 \end{bmatrix}$.

\vspace{0.2cm}\noindent\textbf{Solution 8.}
$AB = \begin{bmatrix} 2 & -1 \\ 3 & 1 \end{bmatrix} \begin{bmatrix} 1 & x \\ y & 2 \end{bmatrix} = \begin{bmatrix} 2-y & 2x-2 \\ 3+y & 3x+2 \end{bmatrix}$.\\
$BA = \begin{bmatrix} 1 & x \\ y & 2 \end{bmatrix} \begin{bmatrix} 2 & -1 \\ 3 & 1 \end{bmatrix} = \begin{bmatrix} 2+3x & -1+x \\ 2y+6 & -y+2 \end{bmatrix}$.\\
Equating entries: (Row 1 Col 1): $2-y = 2+3x \implies y = -3x$.\\
(Row 1 Col 2): $2x-2 = -1+x \implies x = 1$. Consequently, $y = -3(1) = -3$.\\
Verification on Row 2 Col 1: $3+(-3) = 0$; $2(-3)+6 = 0$. Matches. \textbf{$x = 1, y = -3$}.

\vspace{0.2cm}\noindent\textbf{Solution 9.}
$AB = \begin{bmatrix} a+c & b+d \\ c & d \end{bmatrix}$. $BA = \begin{bmatrix} a & a+b \\ c & c+d \end{bmatrix}$.\\
Equating: $a+c = a \implies c = 0$.\\
$b+d = a+b \implies d = a$.\\
Condition: \textbf{$c = 0$} and \textbf{$a = d$}. $b$ can be any real number. 

\vspace{0.2cm}\noindent\textbf{Solution 10.}
$AB = \begin{bmatrix} 8+2 & 4k-6 \\ -4-1 & -2k+3 \end{bmatrix} = \begin{bmatrix} 10 & 4k-6 \\ -5 & -2k+3 \end{bmatrix}$.\\
$BA = \begin{bmatrix} 8-2k & -4+k \\ -2-6 & 2+3 \end{bmatrix} = \begin{bmatrix} 8-2k & k-4 \\ -8 & 5 \end{bmatrix}$.\\
Observe Row 2, Col 1: For $AB$, it is $-5$. For $BA$, it is $-8$. Since $-5 \neq -8$ unconditionally, \textbf{no value of $k$ exists} that allows them to commute.

\vspace{0.2cm}\noindent\textbf{Solution 11.}
$DM = \begin{bmatrix} 2a & 2b \\ 3c & 3d \end{bmatrix}$. $MD = \begin{bmatrix} 2a & 3b \\ 2c & 3d \end{bmatrix}$.\\
Equating implies $2b = 3b \implies b=0$ and $3c = 2c \implies c=0$. $a$ and $d$ are free.\\
$M$ must be a \textbf{diagonal matrix} $\begin{bmatrix} a & 0 \\ 0 & d \end{bmatrix}$.

\vspace{0.2cm}\noindent\textbf{Solution 12.}
We need $Ab_1 = 0$ and $Ab_2 = 0$. Solve $3x_1 + 6x_2 = 0 \implies x_1 = -2x_2$.\\
Let $x_2 = 1 \implies x_1 = -2$. Col 1 is $\begin{bmatrix} -2 \\ 1 \end{bmatrix}$.\\
Let $x_2 = 2 \implies x_1 = -4$. Col 2 is $\begin{bmatrix} -4 \\ 2 \end{bmatrix}$.\\
\textbf{$B = \begin{bmatrix} -2 & -4 \\ 1 & 2 \end{bmatrix}$}. (Other scalar multiples valid).

\vspace{0.2cm}\noindent\textbf{Solution 13.}
Left multiplication $DC=0$ means rows of $D$ dotted with cols of $C$ are 0. Let $D = \begin{bmatrix} x & y \\ z & w \end{bmatrix}$.\\
$DC = \begin{bmatrix} x-2y & -3x+6y \\ z-2w & -3z+6w \end{bmatrix} = \begin{bmatrix} 0 & 0 \\ 0 & 0 \end{bmatrix}$.\\
This requires $x=2y$ and $z=2w$. Let $y=1, w=2$. Then $x=2, z=4$.\\
\textbf{$D = \begin{bmatrix} 2 & 1 \\ 4 & 2 \end{bmatrix}$}.

\vspace{0.2cm}\noindent\textbf{Solution 14.}
$(A - 2B)(A + B) = A(A) + A(B) - 2B(A) - 2B(B) = \mathbf{A^2 + AB - 2BA - 2B^2}$. (Cannot combine middle terms).

\vspace{0.2cm}\noindent\textbf{Solution 15.}
Because the Identity matrix $I$ commutes with any matrix $A$ (i.e., $AI = IA = A$), the algebraic limits of matrices are lifted. $(A+I)^3 = \mathbf{A^3 + 3A^2 + 3A + I}$.

\vspace{0.2cm}\noindent\textbf{Solution 16.}
$u^T v = (1)(4) + (-2)(0) + (3)(-1) = 4 + 0 - 3 = \mathbf{1}$.

\vspace{0.2cm}\noindent\textbf{Solution 17.}
$uv^T = \begin{bmatrix} 1 \\ -2 \\ 3 \end{bmatrix} \begin{bmatrix} 4 & 0 & -1 \end{bmatrix} = \mathbf{\begin{bmatrix} 4 & 0 & -1 \\ -8 & 0 & 2 \\ 12 & 0 & -3 \end{bmatrix}}$.

\vspace{0.2cm}\noindent\textbf{Solution 18.}
$x^T A = \begin{bmatrix} 1 & 2 \end{bmatrix} \begin{bmatrix} 3 & -1 \\ 0 & 4 \end{bmatrix} = \begin{bmatrix} 3 & 7 \end{bmatrix}$.\\
$(x^T A) x = \begin{bmatrix} 3 & 7 \end{bmatrix} \begin{bmatrix} 1 \\ 2 \end{bmatrix} = 3(1) + 7(2) = \mathbf{17}$. The dimension is $1 \times 1$ (a scalar).

\vspace{0.2cm}\noindent\textbf{Solution 19.}
$\text{tr}(M) = m_{11} + m_{22} + m_{33} = 5 + (-3) + 9 = \mathbf{11}$.

\vspace{0.2cm}\noindent\textbf{Solution 20.}
$AB = \begin{bmatrix} 2(1)+(-1)5 & 2(4)+(-1)(-2) \\ 0(1)+3(5) & 0(4)+3(-2) \end{bmatrix} = \begin{bmatrix} -3 & 10 \\ 15 & -6 \end{bmatrix}$. $\text{tr}(AB) = -3 - 6 = \mathbf{-9}$.\\
$BA = \begin{bmatrix} 1(2)+4(0) & 1(-1)+4(3) \\ 5(2)+(-2)0 & 5(-1)+(-2)3 \end{bmatrix} = \begin{bmatrix} 2 & 11 \\ 10 & -11 \end{bmatrix}$. $\text{tr}(BA) = 2 - 11 = \mathbf{-9}$. Property holds.

\vspace{0.2cm}\noindent\textbf{Solution 21.}
$AB = \begin{bmatrix} 1(3)+0(4) & 1(1)+0(-1) \\ -1(3)+2(4) & -1(1)+2(-1) \end{bmatrix} = \begin{bmatrix} 3 & 1 \\ 5 & -3 \end{bmatrix}$. $(AB)^T = \mathbf{\begin{bmatrix} 3 & 5 \\ 1 & -3 \end{bmatrix}}$.\\
$B^T A^T = \begin{bmatrix} 3 & 4 \\ 1 & -1 \end{bmatrix} \begin{bmatrix} 1 & -1 \\ 0 & 2 \end{bmatrix} = \begin{bmatrix} 3+0 & -3+8 \\ 1+0 & -1-2 \end{bmatrix} = \mathbf{\begin{bmatrix} 3 & 5 \\ 1 & -3 \end{bmatrix}}$. They match.

\vspace{0.2cm}\noindent\textbf{Solution 22.}
A matrix $M$ is symmetric if $M^T = M$. Let $M = A^T A$. Take the transpose of the entire expression:
$(A^T A)^T$. Using the transpose reversal rule $(XY)^T = Y^T X^T$, we get:
$(A^T A)^T = (A)^T (A^T)^T = A^T A$. Since the transpose of the product equals the original product, it is rigorously symmetric.

\vspace{0.2cm}\noindent\textbf{Solution 23.}
$uv^T = \begin{bmatrix} u_1 \\ u_2 \end{bmatrix} \begin{bmatrix} v_1 & v_2 \end{bmatrix} = \begin{bmatrix} u_1v_1 & u_1v_2 \\ u_2v_1 & u_2v_2 \end{bmatrix}$.\\
The trace is the sum of the main diagonal: $\text{tr}(uv^T) = u_1v_1 + u_2v_2$.\\
The inner product is $u^T v = \begin{bmatrix} u_1 & u_2 \end{bmatrix} \begin{bmatrix} v_1 \\ v_2 \end{bmatrix} = u_1v_1 + u_2v_2$. Thus, $\text{tr}(uv^T) = u^T v$.

\vspace{0.2cm}\noindent\textbf{Solution 24.}
\textbf{Flaw:} The student assumed matrix multiplication is commutative ($AB = BA$). It is not.\\
\textbf{Correct:} $(A+B)(A+B) = A(A) + A(B) + B(A) + B(B) = \mathbf{A^2 + AB + BA + B^2}$.

\vspace{0.2cm}\noindent\textbf{Solution 25.}
\textbf{Flaw:} The zero-product property applies to real numbers, but not to matrices lacking full rank (non-invertible matrices). \\
\textbf{Correct:} Solve the system $2x_1 + 4x_2 = 0 \implies x_1 = -2x_2$. The correct solution is an infinite set of vectors: $\mathbf{X = c\begin{bmatrix} -2 \\ 1 \end{bmatrix}}$, where $c$ is any real number.

\vspace{0.2cm}\noindent\textbf{Solution 26.}
\textbf{Flaw 1:} Vectors $u$ and $v$ do not have inverses. Inversion is strictly for square matrices.\\
\textbf{Flaw 2:} The outer product $uv^T$ generates a $3 \times 3$ matrix of Rank 1. Since its columns are merely scalar multiples of $u$, the matrix determinant is zero, meaning the matrix $(uv^T)$ itself is completely non-invertible.

\vspace{0.2cm}\noindent\textbf{Solution 27.}
\textbf{Flaw:} The student forgot to reverse the order of the matrices when distributing the transpose. $(AB)^T = B^T A^T$.\\
\textbf{Correct:} $B^T$ is $5 \times 4$. $A^T$ is $4 \times 3$. Multiplying $(5 \times 4) \times (4 \times 3)$ legally yields a \textbf{$5 \times 3$} matrix.

\vspace{0.2cm}\noindent\textbf{Solution 28.}
$w^T \Sigma = \begin{bmatrix} 0.5 & 0.5 \end{bmatrix} \begin{bmatrix} 0.16 & 0.02 \\ 0.02 & 0.25 \end{bmatrix} = \begin{bmatrix} 0.08+0.01 & 0.01+0.125 \end{bmatrix} = \begin{bmatrix} 0.09 & 0.135 \end{bmatrix}$.\\
$w^T \Sigma w = \begin{bmatrix} 0.09 & 0.135 \end{bmatrix} \begin{bmatrix} 0.5 \\ 0.5 \end{bmatrix} = 0.045 + 0.0675 = \mathbf{0.1125}$. (Portfolio Variance).

\vspace{0.2cm}\noindent\textbf{Solution 29.}
$P^2 = \begin{bmatrix} 0.9 & 0.2 \\ 0.1 & 0.8 \end{bmatrix} \begin{bmatrix} 0.9 & 0.2 \\ 0.1 & 0.8 \end{bmatrix} = \begin{bmatrix} 0.81+0.02 & 0.18+0.16 \\ 0.09+0.08 & 0.02+0.64 \end{bmatrix} = \begin{bmatrix} 0.83 & 0.34 \\ 0.17 & 0.66 \end{bmatrix}$.\\
$v_2 = P^2 v_0 = \begin{bmatrix} 0.83 & 0.34 \\ 0.17 & 0.66 \end{bmatrix} \begin{bmatrix} 1 \\ 0 \end{bmatrix} = \mathbf{\begin{bmatrix} 0.83 \\ 0.17 \end{bmatrix}}$. (83\% of users retained in Product A).

\vspace{0.2cm}\noindent\textbf{Solution 30.}
$A^2 = \begin{bmatrix} 0 & 1 & 0 \\ 0 & 0 & 1 \\ 0 & 0 & 0 \end{bmatrix} \begin{bmatrix} 0 & 1 & 0 \\ 0 & 0 & 1 \\ 0 & 0 & 0 \end{bmatrix} = \mathbf{\begin{bmatrix} 0 & 0 & 1 \\ 0 & 0 & 0 \\ 0 & 0 & 0 \end{bmatrix}}$.\\
$A^3 = A^2 A = \begin{bmatrix} 0 & 0 & 1 \\ 0 & 0 & 0 \\ 0 & 0 & 0 \end{bmatrix} \begin{bmatrix} 0 & 1 & 0 \\ 0 & 0 & 1 \\ 0 & 0 & 0 \end{bmatrix} = \mathbf{\begin{bmatrix} 0 & 0 & 0 \\ 0 & 0 & 0 \\ 0 & 0 & 0 \end{bmatrix}}$.\\
This reinforces the \textbf{Zero-Product Trap}. It shows that $A \neq 0$ and $A^2 \neq 0$, yet multiplying them repeatedly can still result in the zero matrix without any scalar terms being inherently zero.

\newpage
% ==========================================
% SECTION 4: CONCEPT CHECKLIST APPENDIX
% ==========================================
\section{Appendix: Concept Checklist Mapping Review}
Use this key to go back and practice specific areas of weakness:

\begin{itemize}
    \item \textbf{Need help with addition, scalars, and size checks?} \\ Review Questions: 1, 2, 3, 4, 5.
    \item \textbf{Need help with standard row-column multiplication?} \\ Review Questions: 6, 7, 14, 15.
    \item \textbf{Need help with $AB=BA$ structural algebraic problems?} \\ Review Questions: 8, 9, 10, 11.
    \item \textbf{Need help finding non-zero matrices that multiply to 0?} \\ Review Questions: 12, 13, 29, 30.
    \item \textbf{Need help with Inner/Outer products and scalar tracking?} \\ Review Questions: 16, 17, 18.
    \item \textbf{Need help with Trace formulas?} \\ Review Questions: 19, 20, 23.
    \item \textbf{Need help with Transposes and matrix symmetry proofs?} \\ Review Questions: 21, 22.
    \item \textbf{Need help recognizing trap answers and conceptual flaws?} \\ Review Questions: 24, 25, 26, 27.
    \item \textbf{Need help applying these tools to real-world matrices?} \\ Review Questions: 28, 29.
\end{itemize}

\end{document}